% ID: SEL-D0013-BA-E01
% ESTADO: cotejada. Segunda pasada visual completa y borrador compilado revisados.
% Fuente: Selectividad/matematicas_aplicadas/matematicas_aplicadas_P_A.pdf
% SHA-256: 0305750f6b6c3cc3edbba71a19cd7cce4276935dff5a0b7d964b02eaedeaa222
% Página: 1; bloque A; ejercicio original 1.

Tres institutos piden presupuesto de alojamiento en Roma en dos agencias de
viajes, que les dan el precio por noche según tipo de habitación: individual,
doble y triple.

La primera agencia ofrece los siguientes precios: individual a 65 euros,
doble a 85 euros y triple a 104 euros. La segunda agencia oferta la individual
a 78 euros, la doble a 83 euros y la triple a 106 euros.

El primer instituto necesita tres habitaciones individuales, quince dobles y
dos triples, el segundo dos individuales, doce dobles y cinco triples y el
tercer instituto una individual, dieciséis dobles y siete triples.
\begin{enumerate}[label=\alph*)]
  \item \textbf{(1 punto)} Exprese, mediante una matriz $A$, los precios de
  las dos agencias según tipo de habitación y con otra matriz $D$ la demanda
  de los tres institutos.
  \item \textbf{(1 punto)} Mediante operaciones con las matrices anteriores,
  calcule el precio por noche que cada agencia facilita a los distintos
  institutos por el total de habitaciones solicitadas. \textquestiondown{}Qué
  agencia le interesaría a cada instituto?
  \item \textbf{(0.5 puntos)} \textquestiondown{}Existe la inversa de la matriz
  $D$? \textquestiondown{}Y de la matriz $A$? Justifique las respuestas.
\end{enumerate}
